\section{Database File Format}
\label{sec:database:format}

The metrics are stored within individual files, one for each metric per day.
The file's consist of three sections, the preamble, the header and the data segments.
All multibyte values are in little-endian format.

The following is an example of the first 128 bytes of a database file:

\begin{table*}[h*)]
    \begin{tabular}{r|rrrrrrrrrrrrrrrr|l}
        \hline
        \multirow{2}{*}{\textbf{offset}}& \multicolumn{16}{c|}{\textbf{Hexadecimal}} & \multirow{2}{*}{\textbf{Text}} \\
         &  \textbf{0}&\textbf{1}&\textbf{2}&\textbf{3}&
            \textbf{4}&\textbf{5}&\textbf{6}&\textbf{7}&
            \textbf{8}&\textbf{9}&\textbf{A}&\textbf{B}&
            \textbf{C}&\textbf{D}&\textbf{E}&\textbf{F}& \\
        \hline
        0000 &
            % PIWEATHERDB
            \cellcolor{yellow}50 & \cellcolor{yellow}49 & \cellcolor{yellow}57 & \cellcolor{yellow}54 &
            \cellcolor{yellow}48 & \cellcolor{yellow}52 & \cellcolor{yellow}44 & \multicolumn{1}{r|}{\cellcolor{yellow}42} &
            %
            \cellcolor{blue!25}30 & \multicolumn{1}{r|}{\cellcolor{blue!25}00} &
            %
            \cellcolor{red!25}01 & \multicolumn{1}{r|}{\cellcolor{red!25}00} &
            %
            \cellcolor{green!25}01 & \multicolumn{1}{r|}{\cellcolor{green!25}00} &
            \cellcolor{grey!10}18 & \cellcolor{grey!10}00 &
            \texttt{PIWTHRDB0.......}\\
        \hline
        0010 & \cellcolor{green!25}0f & \multicolumn{1}{r|}{\cellcolor{green!25}00} &
            \cellcolor{yellow}68 & \cellcolor{yellow}6f & \cellcolor{yellow}6d & \cellcolor{yellow}65 &
            \cellcolor{yellow}2e & \cellcolor{yellow}67 & \cellcolor{yellow}65 & \cellcolor{yellow}69 &
            \cellcolor{yellow}67 & \cellcolor{yellow}65 & \cellcolor{yellow}72 & \cellcolor{yellow}2e & \cellcolor{yellow}63 & \cellcolor{yellow}70 &
            \texttt{..home.geiger.cp}\\
            \cline{2-17}
        0020 &
        \multicolumn{1}{r|}{\cellcolor{yellow}6d} &
            00 & 00 & 00 & 00 & 00 & 00 & 00 & 00 & 00 & 00 & 00 & 00 & 00 & 00 & 00 & \texttt{m...............}\\
        \hline
        0030 &
            % r1 time
            \cellcolor{yellow}80 & \cellcolor{yellow}2a & \cellcolor{yellow}b1 & \cellcolor{yellow}67 &
            \cellcolor{yellow}00 & \cellcolor{yellow}00 & \cellcolor{yellow}00 & \cellcolor{yellow}00 &
            % r1 val
            \cellcolor{green!25}00 & \cellcolor{green!25}00 & \cellcolor{green!25}00 & \cellcolor{green!25}00 &
            \cellcolor{green!25}00 & \cellcolor{green!25}00 & \cellcolor{green!25}31 & \cellcolor{green!25}40 &
            \texttt{.*.g..........1@}\\
            \cline{10-17}
        0040 &
            % r1 unit
            \cellcolor{red!25}23 & \cellcolor{red!25}11 & \cellcolor{red!25}c0 & \cellcolor{red!25}53 &
            \cellcolor{red!25}c6 & \cellcolor{red!25}0b & \cellcolor{red!25}b2 & \multicolumn{1}{r|}{\cellcolor{red!25}05} &
            81 & 2a & b1 & 67 & 00 & 00 & 00 & 00 & \texttt{#..S.....*.g....}\\
        \cline{2-9}
        0050 & 00 & 00 & 00 & 00 & 00 & 00 & 32 & 40 & 23 & 11 & c0 & 53 & c6 & 0b & b2 & 05 & \texttt{......2@#..S....}\\
        \cline{2-17}
        0060 & 82 & 2a & b1 & 67 & 00 & 00 & 00 & 00 & 00 & 00 & 00 & 00 & 00 & 00 & 32 & 40 & \texttt{.*.g..........2@}\\
        \cline{10-17}
        0070 & 23 & 11 & c0 & 53 & c6 & 0b & b2 & \multicolumn{1}{r|}{05} & 83 & 2a & b1 & 67 & 00 & 00 & 00 & 00 & \texttt{#..S.....*.g....}\\
        \cline{2-9}
    \end{tabular}
\end{table*}

\subsection{Preamble}
\label{subsec:database:preamble}

The preamble is the first 16 bytes of the file:

\begin{table*}[h*)]
    \begin{tabular}{r|rrrrrrrrrrrrrrrr|l}
        \hline
        \multirow{2}{*}{\textbf{offset}}& \multicolumn{16}{c|}{\textbf{Hexadecimal}} & \multirow{2}{*}{\textbf{Text}} \\
        &  \textbf{0}&\textbf{1}&\textbf{2}&\textbf{3}&
        \textbf{4}&\textbf{5}&\textbf{6}&\textbf{7}&
        \textbf{8}&\textbf{9}&\textbf{A}&\textbf{B}&
        \textbf{C}&\textbf{D}&\textbf{E}&\textbf{F}& \\
        \hline
        0000 &
        % PIWEATHERDB
        \cellcolor{yellow}50 & \cellcolor{yellow}49 & \cellcolor{yellow}57 & \cellcolor{yellow}54 &
        \cellcolor{yellow}48 & \cellcolor{yellow}52 & \cellcolor{yellow}44 & \multicolumn{1}{r|}{\cellcolor{yellow}42} &
        %
        \cellcolor{blue!25}30 & \multicolumn{1}{r|}{\cellcolor{blue!25}00} &
        %
        \cellcolor{red!25}01 & \multicolumn{1}{r|}{\cellcolor{red!25}00} &
        %
        \cellcolor{green!25}01 & \multicolumn{1}{r|}{\cellcolor{green!25}00} &
        \cellcolor{grey!10}18 & \cellcolor{grey!10}00 &
        \texttt{PIWTHRDB0.......}\\
        \hline
    \end{tabular}
\end{table*}

\begin{table}[h!)]
\begin{tabular}{r l l l l}
\multicolumn{2}{c}{\textbf{offset}} & \textbf{Type} & \textbf{Contents} & \textbf{In example}\\
\hline
00 & 07 & string & \multicolumn{2}{l}{Always ``PIWTHRDB'' which identifies this as a database file} \\
08 & 09 & uint16 & The offset to the start of the main body & 0x0030 \\
0A & 0B & uint16 & The header version & 0x0001 \\
0C & 0D & uint16 & The record version & 0x0001\\
0E & 0F & uint16 & Reserved & \\
\hline
\end{tabular}
\end{table}

The reserved field was originally the record length,
\subsectionpaged{Header}
\label{subsec:database:header}

The header is a variable length block starting from offset 0x0010.
The header block is always padded out with nulls to a multiple of 16\sidenote{this also makes debugging easier as you can use hexdump on the db file}
bytes so that the data block always starts at an even boundary.

Its content is defined by the header version in the preamble.
Currently, the only version defined is 0x0001.

\subsubsection{Header version 0x0001}
\label{subsubsec:database:header:0001}

For version 1 the only value stored in the header is the metric name of the data this file is storing.
\begin{table}[ht]
    \begin{tabular}{r l l l}
        \multicolumn{2}{c}{\textbf{file offset}} & \textbf{Type} & \textbf{Contents} \\
        \hline
        10 & 11 & uint16 & The length of the metric name \\
        12 & n & string & The metric name \\
        \hline
    \end{tabular}
\end{table}

From the earlier example, we have the following for the header:
\begin{table*}[h*)]
    \begin{tabular}{r|rrrrrrrrrrrrrrrr|l}
        \hline
        \multirow{2}{*}{\textbf{offset}}& \multicolumn{16}{c|}{\textbf{Hexadecimal}} & \multirow{2}{*}{\textbf{Text}} \\
        &  \textbf{0}&\textbf{1}&\textbf{2}&\textbf{3}&
        \textbf{4}&\textbf{5}&\textbf{6}&\textbf{7}&
        \textbf{8}&\textbf{9}&\textbf{A}&\textbf{B}&
        \textbf{C}&\textbf{D}&\textbf{E}&\textbf{F}& \\
        \hline
        0010 & \cellcolor{green!25}0f & \multicolumn{1}{r|}{\cellcolor{green!25}00} &
        \cellcolor{yellow}68 & \cellcolor{yellow}6f & \cellcolor{yellow}6d & \cellcolor{yellow}65 &
        \cellcolor{yellow}2e & \cellcolor{yellow}67 & \cellcolor{yellow}65 & \cellcolor{yellow}69 &
        \cellcolor{yellow}67 & \cellcolor{yellow}65 & \cellcolor{yellow}72 & \cellcolor{yellow}2e & \cellcolor{yellow}63 & \cellcolor{yellow}70 &
        \texttt{..home.geiger.cp}\\
        \cline{2-17}
        0020 &
        \multicolumn{1}{r|}{\cellcolor{yellow}6d} &
        00 & 00 & 00 & 00 & 00 & 00 & 00 & 00 & 00 & 00 & 00 & 00 & 00 & 00 & 00 & \texttt{m...............}\\
        \hline
    \end{tabular}
\end{table*}

Here the metric name is 15 bytes long (0x0f) and is set to \Metric{home.geiger.cpm}.
Bytes 0x21\dots0x2f are padding so are unused.

\subsectionpaged{Body}
\label{subsec:database:body}

The body contains the actual recorded data.
It starts from the offset set in the preamble.

Each record is of fixed length and the record format is defined by the record version in the preamble.
Currently, the only version defined is 0x0001.

\subsubsection{Record version 0x0001}
\label{subsubsec:database:record:0001}

For version 1 we store the three critical portions of any reading: the time, the value and the unit of the value.
\begin{table}[ht]
    \begin{tabular}{r l l l}
        \toprule
        \multicolumn{2}{c}{\textbf{file offset}} & \textbf{Type} & \textbf{Contents} \\
        \midrule
        10 & 11 & uint16 & The length of the metric name \\
        12 & n & string & The metric name \\
        \bottomrule
    \end{tabular}
\end{table}

\begin{table*}[h*)]
    \begin{tabular}{r|rrrrrrrrrrrrrrrr|l}
        \hline
        \multirow{2}{*}{\textbf{offset}}& \multicolumn{16}{c|}{\textbf{Hexadecimal}} & \multirow{2}{*}{\textbf{Text}} \\
        &  \textbf{0}&\textbf{1}&\textbf{2}&\textbf{3}&
        \textbf{4}&\textbf{5}&\textbf{6}&\textbf{7}&
        \textbf{8}&\textbf{9}&\textbf{A}&\textbf{B}&
        \textbf{C}&\textbf{D}&\textbf{E}&\textbf{F}& \\
        \hline
        0030 &
        % r1 time
        \cellcolor{yellow}80 & \cellcolor{yellow}2a & \cellcolor{yellow}b1 & \cellcolor{yellow}67 &
        \cellcolor{yellow}00 & \cellcolor{yellow}00 & \cellcolor{yellow}00 & \cellcolor{yellow}00 &
        % r1 val
        \cellcolor{green!25}00 & \cellcolor{green!25}00 & \cellcolor{green!25}00 & \cellcolor{green!25}00 &
        \cellcolor{green!25}00 & \cellcolor{green!25}00 & \cellcolor{green!25}31 & \cellcolor{green!25}40 &
        \texttt{.*.g..........1@}\\
        \cline{10-17}
        0040 &
        % r1 unit
        \cellcolor{red!25}23 & \cellcolor{red!25}11 & \cellcolor{red!25}c0 & \cellcolor{red!25}53 &
        \cellcolor{red!25}c6 & \cellcolor{red!25}0b & \cellcolor{red!25}b2 & \multicolumn{1}{r|}{\cellcolor{red!25}05} &
        81 & 2a & b1 & 67 & 00 & 00 & 00 & 00 & \texttt{#..S.....*.g....}\\
        \cline{2-9}
    \end{tabular}
\end{table*}
